\section{Prolog}
\label{sec:prolog}
\slidesrc{Prolog I: 1--37; Prolog II: 1--9; Prolog IV: 1--15; Prolog V: 1--24; Prolog III: 1--9}

\textbf{Prolog} (PROgramming in LOGic, Colmerauer and Roussel, 1972;
Section~\ref{sec:ai-history}) is the language in which symbolic AI is most directly
\emph{programmed}. A Prolog program is a set of logical statements (facts and rules, i.e.\
Horn clauses of first-order logic, Section~\ref{sec:first-order-logic}); running a program
means asking a question (a \emph{query}) and letting Prolog \emph{prove} it. The professor's
slides are based on material by Peter Schachte (University of Melbourne).

The five decks are presented in \textbf{lecture order}: I (basics), II (lists),
IV (nondeterminism, arithmetic, all-solutions), V (meta-programming and the database),
III (logical foundations: Horn clauses and SLD resolution). Part III comes last, as in class:
it explains \emph{why} everything seen before works.

\begin{nbwhy}[How to try the examples]
\notinslides
The slides show sessions of tuProlog and SICStus Prolog, which answer \code{Yes}/\code{No}.
Today the most common system is \textbf{SWI-Prolog} (free, \url{https://www.swi-prolog.org}):
write the program in a file \code{royals.pl}, start \code{swipl}, load it with
\code{?- [royals].} (or \code{consult(royals)}) and type queries after \code{?-}. SWI-Prolog
answers \code{true}/\code{false} instead of \code{Yes}/\code{No}, names fresh variables
\code{\_123} or \code{\_A} instead of \code{\_G168}, and asks for more answers when you press
\code{;} (or space). In these notes answers are written in SWI-Prolog style, with the slides'
output where it differs.
\end{nbwhy}

\subsection{Prolog basics}
\label{sec:prolog-basics}
\slidesrc{Prolog I: 1--37}

\subsubsection{Predicates and facts}
\slidesrc{Prolog I: 4--5}

In most programming languages the basic abstraction is the \emph{function} (or method): a
mapping from inputs to outputs. In Prolog the basic abstraction is the \textbf{predicate}
(or \textbf{relation}): it indicates \emph{whether or not some entities are related} in a
particular way. In logic programming we may \emph{test} whether entities are related, or
\emph{find} entities that are related to other entities.

\begin{nbexample}{the royal family program (slide 5)}
\begin{minted}{prolog}
parent(elizabeth, charles).
parent(diana, william).
parent(diana, harry).
parent(philip, charles).
parent(charles, william).
parent(charles, harry).
\end{minted}
\begin{itemize}[nosep]
  \item \code{parent} is a \textbf{predicate} (of arity 2, written \code{parent/2});
  \item each line is a separate \textbf{clause}, terminated by a full stop;
  \item each clause is a separate \emph{true statement}, understandable on its own:
        ``Elizabeth is a parent of Charles'';
  \item names of individuals (\code{elizabeth}) are \textbf{atoms}: constants, starting
        with a lower-case letter.
\end{itemize}
A Prolog program may consist of many files, each defining many predicates, each defined by
many clauses.
\end{nbexample}

\begin{center}
\begin{forest}
for tree={draw, rounded corners, font=\small, l sep=6mm, s sep=4mm, edge={->, thick}, fill=nbblue!8}
[{elizabeth $+$ philip}, fill=nbpurple!10
  [charles
    [william, fill=nbgreen!10] [harry, fill=nbgreen!10]]]
\end{forest}
\hspace{1.5cm}
\begin{forest}
for tree={draw, rounded corners, font=\small, l sep=6mm, s sep=4mm, edge={->, thick}, fill=nbblue!8}
[{charles $+$ diana}, fill=nbpurple!10
  [william, fill=nbgreen!10] [harry, fill=nbgreen!10]]
\end{forest}
\captionof{figure}{The family encoded by the \code{parent/2} facts (arrows go from parent to
child).}
\end{center}

\subsubsection{Making queries}
\slidesrc{Prolog I: 6--7}

Prolog systems provide an interactive environment to load and test code. The prompt is
\code{?-}. A source file (extension \code{.pl}) is loaded by writing its name in single
quotes inside square brackets, followed by a full stop: \code{?- ['royals.pl'].} Queries are
typed at the prompt and Prolog answers them:
\begin{minted}[linenos=false]{prolog}
?- parent(diana, harry).
true.
?- parent(elizabeth, diana).
false.
\end{minted}
(The slides' tuProlog prints \code{Yes} and \code{No}.)

\begin{nbwarning}[``false'' means ``not provable'', not ``false in reality'']
\code{parent(elizabeth, diana)} fails because it cannot be derived from the program. Prolog
works under the \textbf{closed-world assumption}: what is not known to be true is treated
as false. Ask \code{?- parent(elizabeth, anne).} and you also get \code{false}, although
Princess Anne is Elizabeth's daughter: she is simply not in the program.
\end{nbwarning}

\subsubsection{Variables and multiple directions}
\slidesrc{Prolog I: 8--11}

\begin{minted}[linenos=false]{prolog}
?- parent(elizabeth, X).
X = charles.
\end{minted}
Read the query as ``find \code{X} such that Elizabeth is a parent of \code{X}''.

\begin{nbdef}{Prolog variables (slide 8)}
Prolog variables begin with a \textbf{capital letter or an underscore}, followed by letters,
digits and underscores (\code{X}, \code{Kid}, \code{\_tmp}). They are very different from
variables in other programming languages: like in Haskell, \textbf{you cannot change the value
of a variable once it is set}. A variable stands for an \emph{unknown} object, as in algebra.
\end{nbdef}

Some queries have more than one answer. Prolog prints one answer and waits: type \code{;}
to ask for the next one, Return to stop.
\begin{minted}[linenos=false]{prolog}
?- parent(diana, X).
X = william ;
X = harry.
\end{minted}
(The slides show \code{X = harry ;} followed by \code{No}: tuProlog only discovers that
there are no more answers after you ask; SWI-Prolog often knows it already and ends with a
full stop.)

\begin{nbfact}{Multiple directions (slides 10--11)}
The same predicate can be queried ``the other way around'':
\begin{minted}[linenos=false]{prolog}
?- parent(X, william).
X = diana ;
X = charles.
\end{minted}
``Find \code{X} such that \code{X} is a parent of William.'' This works because the program
is \textbf{declarative}: it specifies information without specifying how it must be used.
For each argument we may pass a value or a variable; values \emph{constrain} the solutions,
variables do not. \code{?- parent(X, Y).} imposes no constraint and enumerates all six
pairs. Well-written logic programs exhibit this flexibility, although it is not always
practical (see \code{rev/2} in Section~\ref{sec:meta-programming}).
\end{nbfact}

\begin{nbintuition}[A predicate is a table, a query is a database question]
\notinslides
\code{parent/2} is like a two-column database table. \code{parent(diana, X)} is
``\code{SELECT child WHERE parent = 'diana'}''; \code{parent(X, william)} is
``\code{SELECT parent WHERE child = 'william'}''. In Java you would need two methods
(\code{getChildren}, \code{getParents}); in Prolog one relation answers both questions.
\end{nbintuition}

\subsubsection{Compound queries and constraints}
\slidesrc{Prolog I: 12--18}

A query can invoke several predicates separated by commas (``and''). If a variable appears
more than once, all occurrences must have the \emph{same value}, as in algebra.
\begin{minted}[linenos=false]{prolog}
?- parent(elizabeth, X), parent(X, Y).
X = charles, Y = william ;
X = charles, Y = harry.
\end{minted}

\begin{nbexample}{how Prolog solves a compound query (backtracking)}
\notinslides
\begin{enumerate}[nosep]
  \item First goal \code{parent(elizabeth, X)}: scan the facts top-down; the first match is
        \code{parent(elizabeth, charles)}, so \code{X = charles}. Remember this choice point.
  \item Second goal, now \code{parent(charles, Y)}: first match
        \code{parent(charles, william)}, \code{Y = william}. Print the answer.
  \item You type \code{;}: Prolog \textbf{backtracks} to the most recent choice point, the
        second goal, and continues scanning: \code{parent(charles, harry)},
        \code{Y = harry}. Print.
  \item \code{;} again: no more \code{parent(charles, \_)} facts; back to the first goal:
        no other \code{parent(elizabeth, \_)} fact. No more answers.
\end{enumerate}
\end{nbexample}

Let us add gender information (slide 13):
\begin{minted}{prolog}
female(elizabeth).   female(diana).
male(philip).        male(charles).      male(william).     male(harry).
\end{minted}
Predicates can take any number of arguments; the number is the \textbf{arity}. A predicate
with one argument defines a \emph{property} that certain entities may have.

\begin{nbexample}{Motherhood (slides 14--16)}
Give a Prolog query that determines the mother of William.
\begin{minted}[linenos=false]{prolog}
?- parent(X, william), female(X).
X = diana.
?- female(X), parent(X, william).     % equally good
X = diana.
\end{minted}
In the first version Prolog tries \code{X = diana} (fact 2), which is female: answer. (On
\code{;} it would try \code{X = charles}, who is not female: no more answers.) In the
second, it tries \code{X = elizabeth} first (not a parent of William), then
\code{X = diana}: answer.
\end{nbexample}

\begin{nbfact}{Constraints (slides 17--18)}
Think of each part of a compound query as a \textbf{constraint} on the solution: Prolog only
shows solutions that satisfy all of them. Some constraints have no solutions; in general
constraints can be placed in any order, but some orders are more efficient than others, and
some orders may not work as expected (more later).
\begin{minted}[linenos=false]{prolog}
?- parent(X,Y), parent(Y,Z), male(Z), female(X).
X = elizabeth, Y = charles, Z = william ;
X = elizabeth, Y = charles, Z = harry ;
false.
?- parent(X,Y), parent(Y,X).
false.
?- male(X), female(X).
false.
\end{minted}
\end{nbfact}

\subsubsection{Rules: clause bodies}
\slidesrc{Prolog I: 19--21}

\begin{nbdef}{Clause, head, body, fact, rule (slide 19)}
A clause may have two parts, a \textbf{head} and a \textbf{body}, separated by \code{:-}
(read ``\textbf{if}''). The head looks like a predicate with arguments (some may be
variables), the body looks like a query. Meaning: \emph{the head is true if the body is
true}. If the body is empty, \code{:-} is omitted: a clause with no body is a \textbf{unit
clause} or \textbf{fact}. A clause with a body is a \textbf{rule}.
\end{nbdef}

\begin{minted}{prolog}
grandparent(G, C) :-
    parent(G, P),
    parent(P, C).
\end{minted}
Read: ``\code{G} is a grandparent of \code{C} \emph{if} \code{G} is a parent of \code{P}
\emph{and} \code{P} is a parent of \code{C}''. Variables in a rule are universally
quantified: in FOL, $\forall G\,\forall C\,\forall P\,(\mathit{parent}(G,P) \land
\mathit{parent}(P,C) \limp \mathit{grandparent}(G,C))$.

\begin{minted}[linenos=false]{prolog}
?- grandparent(elizabeth, X).
X = william ;
X = harry.
?- grandparent(X, william).
X = elizabeth ;
X = philip ;
false.
\end{minted}

\begin{center}
\begin{forest}
for tree={draw, rounded corners, font=\small\ttfamily, l sep=5mm, s sep=8mm, edge={thick}, align=center}
[{?- grandparent(elizabeth, X)}, fill=nbblue!10
  [{?- parent(elizabeth, P), parent(P, X)}, fill=nbblue!10, edge label={node[midway, right, font=\scriptsize\rmfamily]{rule, G = elizabeth, C = X}}
    [{?- parent(charles, X)}, fill=nbblue!10, edge label={node[midway, right, font=\scriptsize\rmfamily]{P = charles}}
      [{success\\X = william}, fill=nbgreen!15]
      [{success\\X = harry}, fill=nbgreen!15]]]]
\end{forest}
\captionof{figure}{The search tree of \code{?- grandparent(elizabeth, X).}: each node is the
list of goals still to prove; each edge is the choice of a clause. Answers are the leaves
where nothing is left to prove. Section~\ref{sec:logical-foundations} calls this an
\emph{SLD tree}.}
\end{center}

\begin{nbexercise}{P1.1}{Family relationships (slide 21)}
\begin{enumerate}[nosep]
  \item Define a predicate \code{father(Dad, Kid)} that holds when \code{Dad} is the father
        of \code{Kid}.
  \item Define a predicate \code{greatgrandparent(Adult, Child)} that holds when
        \code{Adult} is a great-grandparent of \code{Child}.
\end{enumerate}
\end{nbexercise}
\begin{nbsolution}
\inputminted{prolog}{sections/prolog/prolog-basics/assets/family-solutions.pl}
\code{?- father(X, harry).} gives \code{X = charles} (the parents of Harry are Diana and
Charles; only Charles is male). With the six facts of the program,
\code{greatgrandparent/2} has \emph{no} solutions (the family tree has only three
generations): \code{?- greatgrandparent(A, C).} answers \code{false}. Adding, e.g.,
\code{parent(william, george).} makes \code{?- greatgrandparent(A, george).} answer
\code{A = elizabeth ; A = philip}. \textbf{Lesson}: reuse existing predicates; a
great-grandparent is a parent of a grandparent.
\end{nbsolution}

\subsubsection{Recursive predicates}
\slidesrc{Prolog I: 22--23}

\begin{minted}{prolog}
ancestor(A, A).
ancestor(A, D) :-
    parent(A, C),
    ancestor(C, D).
\end{minted}
The first clause says that everyone is their own ancestor (not the usual definition, but
common in computer science); the second that the parents of your ancestors are your
ancestors. The same code finds both ancestors and descendants:
\begin{minted}[linenos=false]{prolog}
?- ancestor(A, harry).
A = harry ; A = elizabeth ; A = diana ; A = philip ; A = charles ; false.
?- ancestor(elizabeth, D).
D = elizabeth ; D = charles ; D = william ; D = harry ; false.
\end{minted}

\begin{nbexample}{why the answers of \texttt{ancestor(A, harry)} come in that order}
\notinslides
\begin{enumerate}[nosep]
  \item Clause 1 unifies \code{ancestor(A, A)} with \code{ancestor(A, harry)}: \code{A = harry}.
  \item Clause 2: prove \code{parent(A, C), ancestor(C, harry)}, trying the parent facts in order:
  \begin{itemize}[nosep]
    \item \code{parent(elizabeth, charles)}: need \code{ancestor(charles, harry)}; clause 1
          fails (\code{charles} $\neq$ \code{harry}); clause 2: \code{parent(charles, william)}
          then \code{ancestor(william, harry)}: fails (William has no children);
          \code{parent(charles, harry)} then \code{ancestor(harry, harry)}: clause 1 succeeds.
          Answer \code{A = elizabeth}.
    \item \code{parent(diana, william)}: \code{ancestor(william, harry)} fails.
    \item \code{parent(diana, harry)}: \code{ancestor(harry, harry)} succeeds: \code{A = diana}.
    \item \code{parent(philip, charles)}: as for Elizabeth: \code{A = philip}.
    \item \code{parent(charles, william)}: fails; \code{parent(charles, harry)}: \code{A = charles}.
  \end{itemize}
\end{enumerate}
The order of answers follows the order of the facts, exactly as in a depth-first search.
\end{nbexample}

\begin{nbwarning}[Put the base case first and recurse on something smaller]
If the body were written \code{ancestor(C, D), parent(A, C)} (recursive call first), the
query \code{?- ancestor(A, harry).} would find the first answers and then call
\code{ancestor} with both arguments unknown, again and again, without ever consuming a
\code{parent} fact: an infinite loop (\emph{left recursion}). Same logic, different
behaviour: the order of goals matters for termination (Section~\ref{sec:logical-foundations}).
\end{nbwarning}

\subsubsection{Disjunction}
\slidesrc{Prolog I: 24}

Disjunction (``or'') in Prolog uses the \code{;} operator; it is equivalent to writing two
clauses:
\begin{center}
\begin{minipage}{0.42\linewidth}
\begin{minted}[linenos=false]{prolog}
person(X) :-
    ( good(X) ; bad(X) ).
\end{minted}
\end{minipage}
\quad$\equiv$\quad
\begin{minipage}{0.42\linewidth}
\begin{minted}[linenos=false]{prolog}
person(X) :- good(X).
person(X) :- bad(X).
\end{minted}
\end{minipage}
\end{center}

\subsubsection{Terms and variables}
\slidesrc{Prolog I: 25--28}

How can a variable be a data structure? A Prolog variable is \emph{just a term whose value
you don't know yet}. Don't \emph{assign} to variables, \emph{constrain} them. A variable can
be bound to any term, even another variable, but cannot change value once bound. It is
more like a variable in algebra than in a conventional programming language.

\begin{nbdef}{Compound term (slide 26)}
The \textbf{compound term} is Prolog's only data-structuring abstraction. It has a
\textbf{functor}, which is an atom, and one or more \textbf{arguments}, which can be any
terms. Syntax: functor first, then the arguments in parentheses, separated by commas. The
number of arguments is the \textbf{arity}. No need to declare it: just use it.
\end{nbdef}

Compound terms look just like predicate invocations (Whiteboard~\ref{wb:terms-atoms}: the
same syntax builds terms and atoms), but they are \emph{data}. Use them where you would use
(a pointer to) a \code{struct} in C; the functor indicates the kind of term:
\begin{center}
\code{position(20, 10)} \qquad \code{employee(1234, 'Jones', 'James', 100000)}
\end{center}
No allocation, no assignment to members: just write it down. Terms may contain variables,
whether or not their values are determined: \code{position(X, Y)},
\code{employee(Num, Surname, Firstname, Salary)}.

\begin{center}
\small
\begin{tabular}{@{}lll@{}}
\toprule
\textbf{Kind of term} & \textbf{Examples} & \textbf{Rule} \\
\midrule
atom & \code{elizabeth}, \code{[]}, \code{'Jones'}, \code{+} & lower-case start, or quoted \\
number & \code{42}, \code{-3}, \code{3.14} & \\
variable & \code{X}, \code{Kid}, \code{\_}, \code{\_tmp} & upper-case or \code{\_} start \\
compound & \code{position(20,10)}, \code{s(s(0))}, \code{[1,2]} & functor(args) \\
\bottomrule
\end{tabular}
\captionof{table}{The kinds of Prolog terms (added). Atoms and numbers are called
\emph{atomic}; a term with no variables is \emph{ground}.}
\end{center}

\subsubsection{Unification}
\slidesrc{Prolog I: 29--33}

In Prolog, \code{=} is used both for giving values to variables and for comparing terms:
\begin{minted}[linenos=false]{prolog}
?- 4 = 4.            true.
?- 4 = 5.            false.
?- X = 5.            X = 5.
?- 5 = V.            V = 5.
?- X = Y.            X = Y.
\end{minted}
(For the last query the slides show \code{X = \_G168, Y = \_G168}: both variables are bound
to the same fresh variable.) Variables appearing in different queries are unrelated.

\begin{nbdef}{Unification (slide 30)}
\code{=} performs \textbf{unification}: it tries to make two terms the \emph{same} by binding
the variables appearing in them. It is a powerful sort of pattern matching, and Prolog also
uses it for argument passing (calling \code{parent(diana, X)} unifies the call with the
head of each clause).
\begin{minted}[linenos=false]{prolog}
?- foo(A,2,C,D) = foo(1,X,Y,X).
A = 1, X = 2, D = 2, C = Y.
?- bar(X, X) = bar(1, 2).
false.
\end{minted}
\end{nbdef}

\begin{nbfact}{How to unify two terms (the algorithm, added)}
\notinslides
To unify $s$ and $t$:
\begin{enumerate}[nosep]
  \item if one of them is an unbound variable, bind it to the other (if both are the same
        variable, do nothing);
  \item if both are atomic (atoms or numbers), they unify iff they are identical;
  \item if both are compound, they unify iff they have the \textbf{same functor and arity}
        and their arguments unify \textbf{pairwise}, left to right, carrying the bindings
        along;
  \item otherwise they do not unify.
\end{enumerate}
The result is the \textbf{most general unifier} (mgu): the bindings that make the terms
equal while committing to as little as possible.
\end{nbfact}

\begin{nbexample}{\texttt{foo(A,2,C,D) = foo(1,X,Y,X)} step by step}
Same functor \code{foo}, same arity 4: unify the arguments pairwise.
\begin{center}
\small
\begin{tabular}{@{}clll@{}}
\toprule
\textbf{Pair} & \textbf{Left} & \textbf{Right} & \textbf{Result} \\
\midrule
1 & \code{A} & \code{1} & bind \code{A = 1} \\
2 & \code{2} & \code{X} & bind \code{X = 2} \\
3 & \code{C} & \code{Y} & bind \code{C = Y} (two variables, now the same unknown) \\
4 & \code{D} & \code{X} & \code{X} is already \code{2}: bind \code{D = 2} \\
\bottomrule
\end{tabular}
\end{center}
For \code{bar(X, X) = bar(1, 2)}: pair 1 binds \code{X = 1}; pair 2 must unify
\code{1} with \code{2}: different atomic terms, failure.
\end{nbexample}

\begin{nbexercise}{P1.2}{Unification (slides 31--33)}
Find bindings that make these pairs of terms the same:
(1) \code{f(X, 17) = f(Y, Y)}; \quad (2) \code{f(42, U, a) = f(V, f(V,W), W)}.
(\code{f} is used with arity 2 and 3: no problem, they are different functors.)
\end{nbexercise}
\begin{nbsolution}
(1) Pair 1: \code{X = Y}. Pair 2: \code{17 = Y}, so \code{Y = 17}, hence \code{X = 17}.
Answer: \code{X = 17, Y = 17}.

(2) Pair 1: \code{42 = V}, so \code{V = 42}. Pair 2: \code{U = f(V, W)} $=$ \code{f(42, W)}.
Pair 3: \code{a = W}, so \code{W = a}, and therefore \code{U = f(42, a)}. Answer:
\code{V = 42, W = a, U = f(42, a)}.
\end{nbsolution}

\begin{nbwarning}[Unification is not assignment, and not arithmetic]
\code{X = 1 + 2} binds \code{X} to the \emph{term} \code{+(1,2)}, not to \code{3}
(arithmetic needs \code{is}, Section~\ref{sec:arithmetic-nondeterminism}). And
\code{X = f(X)}: standard unification would need the \emph{occurs check} (a variable cannot
be bound to a term containing itself) and fail; SWI-Prolog skips the check for speed and
builds a cyclic term. Avoid such goals.
\end{nbwarning}

\subsubsection{Recursive structures: numbers as terms}
\slidesrc{Prolog I: 34--37}

Any argument of a compound term can be any term, including a term with the same functor:
this is how Prolog implements recursive structures (lists, trees, \ldots). A good strategy to
write predicates on recursive structures:
\begin{enumerate}[nosep]
  \item write a predicate that \emph{recognises} the recursive data structure;
  \item use it as a ``skeleton'' for other predicates, adding arguments and goals.
\end{enumerate}

For example, represent a natural number as either \code{0} or \code{s(X)} (``one more than
\code{X}'', the \emph{successor}), where \code{X} is a natural number: $3$ is
\code{s(s(s(0)))}. (This is \emph{not} how Prolog really does maths.)
\begin{minted}{prolog}
nat(0).                     % 0 is a natural number
nat(s(X)) :- nat(X).        % s(X) is a nat if X is

add(0, Y, Y).               % adding 0 to Y gives Y
add(s(X), Y, s(Z)) :-       % adding X+1 to Y gives Z+1
    add(X, Y, Z).           %   where Z is X + Y
\end{minted}
The structure of \code{add/3} follows that of \code{nat/1}. We refer to predicates as
\code{name/arity} because Prolog allows different predicates with the same name and
different arity.

\begin{nbexample}{\texttt{add(s(s(0)), s(s(s(0))), X)}: $2 + 3$}
\begin{align*}
&\code{add(s(s(0)), s(s(s(0))), X)} && \reason{clause 2: \code{X = s(Z1)}} \\
&\code{add(s(0), s(s(s(0))), Z1)} && \reason{clause 2: \code{Z1 = s(Z2)}} \\
&\code{add(0, s(s(s(0))), Z2)} && \reason{clause 1: \code{Z2 = s(s(s(0)))}}
\end{align*}
Back-substituting: \code{X = s(Z1) = s(s(Z2)) = s(s(s(s(s(0)))))}, i.e.\ $5$. Check:
$2 + 3 = 5$.
\end{nbexample}

\begin{nbfact}{Example queries (slide 36)}
\begin{minted}[linenos=false, fontsize=\footnotesize]{prolog}
?- nat(s(s(0))).                         true.
?- nat(X).                               X = 0 ; X = s(0) ; X = s(s(0)) ; ...
?- add(s(s(0)), X, Y).                   Y = s(s(X)).
?- add(s(s(0)), X, s(s(s(s(s(0)))))).    X = s(s(s(0))).        % 5 - 2
?- add(X, s(s(s(0))), s(s(s(s(s(0)))))). X = s(s(0)) ;  false.  % 5 - 3
?- add(X, Y, s(s(0))).                   X = 0, Y = s(s(0)) ;
                                         X = s(0), Y = s(0) ;
                                         X = s(s(0)), Y = 0 ; false.
\end{minted}
The same definition of addition also computes \textbf{subtraction} and enumerates all ways
of writing a number as a sum: multiple directions again. \code{nat(X)} generates the natural
numbers forever.
\end{nbfact}

\begin{nbexercise}{P1.3}{Multiplication (slide 37)}
Define \code{mult(X, Y, Z)}, true when \code{X} times \code{Y} equals \code{Z}, using
\code{add/3}.
\end{nbexercise}
\begin{nbsolution}
\inputminted{prolog}{sections/prolog/prolog-basics/assets/mult-solution.pl}
This follows the skeleton of \code{nat/1} on the first argument: $0 \cdot Y = 0$ and
$(X+1)\cdot Y = X\cdot Y + Y$. Trace of $2 \cdot 3$: \code{mult(s(s(0)), s(s(s(0))), Z)}
calls \code{mult(s(0), \ldots, W1)}, which calls \code{mult(0, \ldots, W0)}, giving
\code{W0 = 0}; then \code{add(0, s(s(s(0))), W1)} gives $W_1 = 3$; then
\code{add(s(s(s(0))), s(s(s(0))), Z)} gives $Z = 6$ = \code{s(s(s(s(s(s(0))))))}.
\end{nbsolution}

\subsection{Lists}
\label{sec:prolog-lists}
\slidesrc{Prolog II: 1--9}

\subsubsection{List syntax}
\slidesrc{Prolog II: 5}

Lists are very widely used in Prolog, and have a special syntax:
\begin{itemize}[nosep]
  \item \code{[]} is the \textbf{empty list};
  \item \code{[H|T]} is the list with \textbf{head} \code{H} (the first element) and
        \textbf{tail} \code{T} (the list of the remaining elements);
  \item \code{[1,2]} is a list of two elements; \code{[A,B,C|D]} is a list whose first three
        elements are \code{A}, \code{B}, \code{C} and whose remaining elements form the
        list \code{D}.
\end{itemize}
A recogniser, in the ``skeleton'' style of Section~\ref{sec:prolog-basics}:
\begin{minted}{prolog}
list([]).
list([_|Tail]) :- list(Tail).
\end{minted}

\begin{nbfact}{Lists are ordinary terms}
The bracket notation is special syntax, but lists are ordinary compound terms: the list
functor is the dot \code{'.'} with arity 2 (in SWI-Prolog~7 and later it is called
\code{'[|]'}). So
\[
\code{[1,2,3]} \;=\; \code{[1|[2,3]]} \;=\; \code{[1|[2|[3|[]]]]} \;=\;
\code{'.'(1, '.'(2, '.'(3, [])))}
\]
\end{nbfact}

\begin{center}
\begin{forest}
for tree={draw, rounded corners, font=\small\ttfamily, l sep=4mm, s sep=8mm, edge={thick}}
[{'.'}, fill=nbpurple!12 [1, fill=nborange!12] [{'.'}, fill=nbpurple!12 [2, fill=nborange!12] [{'.'}, fill=nbpurple!12 [3, fill=nborange!12] [{[\,]}, fill=nbgray!15]]]]
\end{forest}
\captionof{figure}{The list \code{[1,2,3]} as a tree: a chain of dot cells, each holding one
element (left) and the rest of the list (right), ending with \code{[]}. Exactly a linked
list.}
\end{center}

\begin{nbexample}{unifying lists}
\notinslides
\begin{center}
\small
\begin{tabular}{@{}ll@{}}
\toprule
\textbf{Query} & \textbf{Answer} \\
\midrule
\code{[H|T] = [a,b,c]} & \code{H = a, T = [b,c]} \\
\code{[H|T] = [a]} & \code{H = a, T = []} \\
\code{[H|T] = []} & \code{false} (the empty list has no head) \\
\code{[A,B|R] = [1,2,3,4]} & \code{A = 1, B = 2, R = [3,4]} \\
\code{[A,B] = [1,2,3]} & \code{false} (different lengths) \\
\bottomrule
\end{tabular}
\end{center}
\end{nbexample}

\subsubsection{List membership}
\slidesrc{Prolog II: 6--7}

\code{member/2} checks list membership. It is a built-in of SWI-Prolog, but could be
defined as:
\begin{minted}{prolog}
member(X, [X|_]).        % X is member of a list beginning with X
member(X, [_|L]) :-      % X is member of a list
    member(X, L).        %   if it is member of its tail
\end{minted}

\begin{nbfact}{Member in action (slide 7)}
\begin{minted}[linenos=false, fontsize=\footnotesize]{prolog}
?- member(c, [a,b,c,d]).          true.
?- member(e, [a,b,c,d]).          false.
?- member(X, [a,b,c]).            X = a ; X = b ; X = c.
?- L = [_,_,_], member(a, L).     L = [a, _, _] ; L = [_, a, _] ; L = [_, _, a].
\end{minted}
The same predicate \emph{checks} membership, \emph{enumerates} the elements, and even
\emph{builds} lists containing a given element.
\end{nbfact}

\begin{nbexample}{trace of \texttt{member(c, [a,b,c,d])}}
\notinslides
\begin{center}
\small
\begin{tabular}{@{}lll@{}}
\toprule
\textbf{Goal} & \textbf{Clause 1 \code{member(X,[X|\_])}} & \textbf{Clause 2} \\
\midrule
\code{member(c, [a,b,c,d])} & \code{c} vs \code{a}: fail & recurse on \code{[b,c,d]} \\
\code{member(c, [b,c,d])} & \code{c} vs \code{b}: fail & recurse on \code{[c,d]} \\
\code{member(c, [c,d])} & \code{c} vs \code{c}: \textbf{success} & \\
\bottomrule
\end{tabular}
\end{center}
For \code{member(e, [a,b,c,d])} the recursion reaches \code{member(e, [])}: neither clause
unifies with an empty list, so the query fails.
\end{nbexample}

\subsubsection{List concatenation: append}
\slidesrc{Prolog II: 8--9}

\code{append/3} concatenates two lists (\code{++} in Haskell). Again a built-in, but could be
defined as:
\begin{minted}{prolog}
append([], L, L).              % [] ++ L gives L
append([J|K], L, [J|KL]) :-    % to ++ a list with first element J and remainder K
    append(K, L, KL).          %   to list L, determine K ++ L and then add J to front
\end{minted}
Compare with Haskell: \code{append [] l = l} and \code{append (j:k) l = j : append k l}.
The same, but Haskell can only append lists, while Prolog's \code{append} is more flexible:

\begin{nbfact}{Append in action (slide 9)}
\begin{minted}[linenos=false, fontsize=\footnotesize]{prolog}
?- append([a,b,c], [d,e], L).      L = [a,b,c,d,e].
?- append(L, [d,e], [a,b,c,d,e]).  L = [a,b,c] ;  false.
?- append([a,b,c], L, [a,b,c,d,e]).L = [d,e].
?- append(K, L, [a,b]).            K = [], L = [a,b] ;
                                   K = [a], L = [b] ;
                                   K = [a,b], L = [] ; false.
\end{minted}
Concatenation, ``what must I prepend?'', ``what remains?'', and ``all ways to split a
list'': four functions in one relation.
\end{nbfact}

\begin{nbexample}{\texttt{append([a,b], [c], L)} step by step}
\begin{align*}
&\code{append([a,b], [c], L)} && \reason{clause 2: \code{J=a, K=[b], L=[a|KL1]}} \\
&\code{append([b], [c], KL1)} && \reason{clause 2: \code{KL1=[b|KL2]}} \\
&\code{append([], [c], KL2)} && \reason{clause 1: \code{KL2=[c]}}
\end{align*}
So \code{L = [a|[b|[c]]] = [a,b,c]}. The work is proportional to the length of the
\emph{first} list.
\end{nbexample}

\begin{center}
\begin{forest}
for tree={draw, rounded corners, font=\small\ttfamily, l sep=5mm, s sep=5mm, edge={thick}, align=center}
[{append(K, L, [a,b])}, fill=nbblue!10
  [{success\\K=[], L=[a,b]}, fill=nbgreen!15, edge label={node[midway, left, font=\scriptsize\rmfamily]{cl.\ 1}}]
  [{append(K1, L, [b])\\K=[a|K1]}, fill=nbblue!10, edge label={node[midway, right, font=\scriptsize\rmfamily]{cl.\ 2}}
    [{success\\K=[a], L=[b]}, fill=nbgreen!15, edge label={node[midway, left, font=\scriptsize\rmfamily]{cl.\ 1}}]
    [{append(K2, L, [])\\K1=[b|K2]}, fill=nbblue!10, edge label={node[midway, right, font=\scriptsize\rmfamily]{cl.\ 2}}
      [{success\\K=[a,b], L=[]}, fill=nbgreen!15, edge label={node[midway, left, font=\scriptsize\rmfamily]{cl.\ 1}}]
      [{fail\\{[]} has no head}, fill=nbred!12, edge label={node[midway, right, font=\scriptsize\rmfamily]{cl.\ 2}}]]]]
\end{forest}
\captionof{figure}{Search tree of \code{?- append(K, L, [a,b]).}: Prolog explores it depth-first,
left to right; each green leaf is one answer, in the order shown in the slides; the red leaf
is why the last \code{;} answers \code{false}.}
\end{center}

\begin{nbwhy}
\notinslides
Lists are Prolog's workhorse: an agent's list of pending tasks, a path in a route planner,
the moves of a game (the exam exercises often ask to move elements between lists with
\code{select/3}, \code{append/3}, \code{member/2}). \code{append/3} used ``backwards'' is
also a compact way to express search: ``find a split of the list such that\ldots''.
\end{nbwhy}

\subsection{Nondeterminism, arithmetic, all solutions}
\label{sec:arithmetic-nondeterminism}
\slidesrc{Prolog IV: 1--15}

\subsubsection{Arithmetic in Prolog: a short primer}
\notinslides

The Prolog~IV deck uses arithmetic without introducing it (that part was in a deck not
shown in class). The essentials:

\begin{center}
\small
\renewcommand{\arraystretch}{1.15}
\begin{tabularx}{\linewidth}{@{}l l X@{}}
\toprule
\textbf{Operator} & \textbf{Example} & \textbf{Meaning} \\
\midrule
\code{X is Expr} & \code{X is 2 + 3} gives \code{X = 5} & \textbf{evaluate} the right side (must be fully known), unify the result with the left side \\
\code{=:=}, \code{=\textbackslash=} & \code{1 + 2 =:= 3} true & arithmetic equal / not equal (both sides evaluated) \\
\code{<}, \code{>}, \code{=<}, \code{>=} & \code{X >= 0} & arithmetic comparison (both sides must be known numbers) \\
\code{=} & \code{X = 2 + 3} gives \code{X = 2+3} & \textbf{unification}: no evaluation, \code{X} becomes the term \code{+(2,3)} \\
\code{==} & \code{X == Y} & identical terms, without binding (Section~\ref{sec:meta-programming}) \\
\bottomrule
\end{tabularx}
\captionof{table}{The four ``equalities'' of Prolog. Note \code{=<} (not \code{<=}).}
\end{center}

\begin{nbwarning}[The classic \texttt{=} vs.\ \texttt{is} trap]
\begin{itemize}[nosep]
  \item \code{X = 3 + 4, X == 7} fails: \code{X} is the term \code{3+4}, not the number 7.
  \item \code{7 is X + 4} with \code{X} unknown raises an error (``Arguments are not
        sufficiently instantiated''): \code{is} cannot solve equations, it only evaluates.
        Arithmetic predicates therefore work in \emph{one direction only}, unlike
        \code{add/3} on \code{s(...)} numbers.
  \item \code{N is N + 1} always fails (for a bound \code{N}): variables cannot change value.
        Use a new variable: \code{N1 is N + 1}.
\end{itemize}
\end{nbwarning}

\subsubsection{Managing nondeterminism}
\slidesrc{Prolog IV: 4--7}

A common Prolog programming error is treating multiple clauses as if they formed an
\textbf{if-then-else} construct. This is an \emph{erroneous} definition of absolute value:
\begin{minted}{prolog}
absolute_value(X, X) :- X >= 0.
absolute_value(X, Y) :- Y is -X.
\end{minted}
\begin{minted}[linenos=false]{prolog}
?- absolute_value(-3, V).      V = 3.
?- absolute_value(42, V).      V = 42 ;
                               V = -42.         % wrong!
\end{minted}
The problem: the second clause promises that $|n| = -n$ \emph{for any} $n$, regardless of its
sign. Each clause must be \emph{true on its own}, because Prolog may use any of them on
backtracking. The correct definition:
\begin{minted}{prolog}
absolute_value(X, X) :- X >= 0.
absolute_value(X, Y) :- X < 0, Y is -X.
\end{minted}
Now \code{?- absolute_value(42, V).} gives only \code{V = 42} (the second clause fails at
\code{42 < 0}).

\begin{nbintuition}[Read every clause as a separate true statement]
In a procedural language, ``else'' means ``in all other cases''. Prolog has no implicit
``else'': clause 2 is a logical statement ``for all \code{X}, the absolute value of \code{X}
is \code{-X}'', which is false for positive numbers. A good test: read each clause aloud
alone. If it is not true, add the missing condition (the \emph{guard}).
\end{nbintuition}

\begin{nbexercise}{P4.1}{Maximum of two numbers (slides 6--7)}
Define \code{maximum(X, Y, Z)} such that \code{Z} is the larger of \code{X} and \code{Y}.
\end{nbexercise}
\begin{nbsolution}
Two clauses, each with its own guard, so each is true on its own:
\code{maximum(X, Y, X) :- X >= Y.} \quad and \quad \code{maximum(X, Y, Y) :- X < Y.}
\\ \code{?- maximum(3, 7, Z).} gives \code{Z = 7}; \code{?- maximum(7, 3, Z).} gives
\code{Z = 7} and nothing else on backtracking. The guards are mutually exclusive
($\geq$ vs.\ $<$), so exactly one clause applies for every pair of numbers.
\end{nbsolution}

\subsubsection{Fibonacci numbers}
\slidesrc{Prolog IV: 8}

Fibonacci numbers: $1, 1, 2, 3, 5, 8, 13, \ldots$ (each is the sum of the previous two).
\begin{minted}{prolog}
fib(0, 1).
fib(1, 1).
fib(N, F) :-
    N1 is N - 1,
    N2 is N - 2,
    fib(N1, F1),
    fib(N2, F2),
    F is F1 + F2.
\end{minted}
\code{?- fib(4, F).} gives \code{F = 5}. The slide asks three questions:

\begin{enumerate}
  \item \textbf{What if we change the order of goals? What if we do \code{F is F1 + F2}
        first?} Then \code{is} is called while \code{F1} and \code{F2} are still unknown:
        error ``Arguments are not sufficiently instantiated''. Arithmetic goals must come
        \emph{after} the goals that compute their inputs. (Moving \code{N1 is N - 1} after
        \code{fib(N1, F1)} breaks it in the same way, and worse: \code{fib(N1, F1)} with
        unknown \code{N1} would match the facts and then loop.)
  \item \textbf{Each call makes two recursive calls. Are there better ways?} The number of
        calls $T(n)$ satisfies $T(n) = 1 + T(n-1) + T(n-2)$ with $T(0) = T(1) = 1$:
        \begin{center}
        \small
        \begin{tabular}{@{}l*{7}{r}@{}}
        \toprule
        $n$ & 2 & 3 & 4 & 5 & 10 & 20 & 30 \\
        \midrule
        $\mathit{fib}(n)$ & 2 & 3 & 5 & 8 & 89 & 10\,946 & 1\,346\,269 \\
        calls $T(n)$ & 3 & 5 & 9 & 15 & 177 & 21\,891 & 2\,692\,537 \\
        \bottomrule
        \end{tabular}
        \end{center}
        Exponential: the same values are recomputed again and again. Better: carry the last
        two values as accumulators (linear time), or remember computed values with
        \code{asserta} (Section~\ref{sec:meta-programming}).
\end{enumerate}

\begin{nbwarning}[Hidden bug: asking for more answers loops forever]
\notinslides
After \code{F = 5}, type \code{;}. Prolog backtracks into \code{fib(1, F1)} and tries clause 3
\emph{also for N = 1} (nothing forbids it): \code{N1 = 0}, \code{N2 = -1}, then
\code{fib(-1, \_)}, \code{fib(-2, \_)}, \ldots{} forever (stack overflow). Same disease as the
absolute value: clause 3 is only true for $N > 1$. Fix: \code{fib(N, F) :- N > 1, N1 is N -
1, \ldots}
\end{nbwarning}

\subsubsection{Lists and arithmetic}
\slidesrc{Prolog IV: 9--13}

Summing a list combines arithmetic with list processing:
\begin{minted}{prolog}
sumlist([], 0).                 % sum([]) = 0
sumlist([X|Xs], Sum) :-         % sum([X|Xs]) = Sum if
    sumlist(Xs, Sum1),          %   sum(Xs) = Sum1 and
    Sum is Sum1 + X.            %   Sum = Sum1 + X
\end{minted}
This can be done more efficiently with an \textbf{accumulator}:
\begin{minted}{prolog}
sumlist(L, Sum) :-              % sumlist/2 defined in terms of sumlist/3
    sumlist(L, 0, Sum).         %   (third argument is the accumulator)
sumlist([], Sum, Sum).          % sum([]) + Sum = Sum
sumlist([X|Xs], Sum0, Sum) :-   % sum([X|Xs]) + Sum0 = Sum if
    Sum1 is Sum0 + X,           %   Sum1 = Sum0 + X, and
    sumlist(Xs, Sum1, Sum).     %   sum(Xs) + Sum1 = Sum
\end{minted}

\begin{nbexample}{the two versions on \texttt{[1,2,3]}}
\begin{center}
\small
\begin{tabular}{@{}l|l@{}}
\textbf{Naive} & \textbf{Accumulator} \\
\hline
\code{sumlist([1,2,3], S)} & \code{sumlist([1,2,3], 0, S)} \\
\quad \code{sumlist([2,3], S1)} & \quad \code{sumlist([2,3], 1, S)} \\
\quad\quad \code{sumlist([3], S2)} & \quad\quad \code{sumlist([3], 3, S)} \\
\quad\quad\quad \code{sumlist([], S3)}: \code{S3 = 0} & \quad\quad\quad \code{sumlist([], 6, S)}: \code{S = 6} \\
\quad\quad \code{S2 is 0 + 3} $= 3$ & \\
\quad \code{S1 is 3 + 2} $= 5$ & \\
\code{S is 5 + 1} $= 6$ & \\
\end{tabular}
\end{center}
The naive version must keep every call open waiting for the recursive result, then adds on
the way back; the accumulator version adds on the way \emph{down}, and the recursive call is
the \emph{last} goal (\textbf{tail recursion}): Prolog can reuse the same stack frame, like a
loop. Both give \code{S = 6}.
\end{nbexample}

\paragraph{List maximum (slide 10).}
\begin{minted}{prolog}
maxlist([X], X).
maxlist([X|Xs], Max) :-
    maxlist(Xs, Max1),
    maximum(Max1, X, Max).
\end{minted}
No clause for \code{[]}: the base case is the singleton list, because \code{[]} has no
maximum. Example: \code{maxlist([3,1,4], M)}: \code{maxlist([4], M2)} gives 4, then
\code{maximum(4, 1, M1)} gives 4, then \code{maximum(4, 3, M)} gives \code{M = 4}.

\paragraph{List length (slide 11).} \code{length/2} relates a list to its length:
\begin{minted}[linenos=false]{prolog}
?- length([a,b,c], 3).     true.
?- length([a,b,c], N).     N = 3.
?- length(L, 3).           L = [_, _, _].
?- length(L, N).           L = [], N = 0 ;  L = [_], N = 1 ;  L = [_, _], N = 2 ; ...
\end{minted}
It is built in most Prolog systems; in SICStus it is in a library, loaded with
\code{:- use\_module(library(lists)).}

\begin{nbfact}{List handling with \texttt{append} and \texttt{length} (slides 12--13)}
\begin{minted}{prolog}
member(X, L) :-             % X is a member of L if L is of the form
    append(_, [X|_], L).    %   any list ++ [X] ++ any list

element(N, L, E) :-         % E is the Nth element of L (counting from 0)
    append(F, [E|_], L),    %   if L = F ++ [E] ++ any list
    length(F, N).           %   and the length of F is N

take(N, L, T) :-            % T is the first N elements of L
    append(T, _, L),        %   if L = T ++ any list
    length(T, N).           %   and the length of T is N

drop(N, L, T) :-            % T is what is left after removing the first N elements
    append(F, T, L),        %   if L = F ++ T
    length(F, N).           %   and the length of F is N

listof([], _).              % [] is a list of anything
listof([E|Es], E) :-        % [E|Es] is a list of Es if
    listof(Es, E).          %   Es is
\end{minted}
\code{?- element(1, [a,b,c], E).} gives \code{E = b}; \code{?- take(2, [a,b,c], T).} gives
\code{T = [a,b]}; \code{?- drop(1, [a,b,c], T).} gives \code{T = [b,c]};
\code{?- listof([x,x,x], x).} is true.
\end{nbfact}

\begin{nbwarning}[Comments of slide 13]
The comments of \code{drop/3} on the slide are copied from \code{take/3} (``T is the first N
elements of L'', ``if L = T ++ any list''): the correct reading is the one above, where
\code{F} (not \code{T}) is the prefix of length \code{N}.
\end{nbwarning}

\begin{nbintuition}[Generate and test]
\code{element}, \code{take} and \code{drop} all follow the pattern ``\code{append}
\emph{generates} all the ways to split the list; \code{length} \emph{tests} which split has
the right size''. Simple and declarative, though not the fastest (a direct recursion stops
earlier).
\end{nbintuition}

\subsubsection{All-solutions predicates}
\slidesrc{Prolog IV: 14--15}

Normally Prolog gives one answer at a time. Built-in predicates collect them all:

\begin{nbdef}{findall, setof, bagof (slide 15)}
\begin{itemize}[nosep]
  \item \code{findall(T, G, L)}: \code{L} is the list of all terms \code{T} that satisfy goal
        \code{G}, in the order solutions are found; variables in \code{G} are left
        uninstantiated; deterministic (if there are no solutions, \code{L = []}).
  \item \code{setof(T, G, L)}: \code{L} is a \textbf{non-empty} list of all terms \code{T}
        that satisfy \code{G}, \textbf{sorted with duplicates removed}; variables appearing
        only in \code{G} can be instantiated; can be nondeterministic (one answer for each
        different instantiation of those variables).
  \item \code{bagof(T, G, L)}: like \code{setof/3} but without the sorting.
\end{itemize}
The template term \code{T} is left uninstantiated by all of these predicates.
\end{nbdef}

\begin{nbexample}{all solutions on the royal family}
\notinslides
\begin{minted}[linenos=false, fontsize=\footnotesize]{prolog}
?- findall(C, parent(charles, C), L).    L = [william, harry].
?- findall(C, parent(P, C), L).          L = [charles, william, harry, charles, william, harry].
?- findall(C, parent(harry, C), L).      L = [].
?- bagof(C, parent(harry, C), L).        false.
?- bagof(C, parent(P, C), L).            P = charles,   L = [william, harry] ;
                                         P = diana,     L = [william, harry] ;
                                         P = elizabeth, L = [charles] ;
                                         P = philip,    L = [charles].
?- setof(C, P^parent(P, C), L).          L = [charles, harry, william].
\end{minted}
\code{findall} ignores \code{P}; \code{bagof} groups the answers by each value of the free
variable \code{P}; \code{P\textasciicircum{}} (``there exists a \code{P}'') tells
\code{setof}/\code{bagof} not to group by \code{P}. \code{setof} sorts and removes duplicates.
\end{nbexample}

\begin{nbwhy}
\notinslides
All-solutions predicates turn a \emph{search} into a \emph{data structure}: ``the list of all
free rooms'', ``all plans applicable now''. Agent interpreters do exactly this when they
compute the set of relevant plans for an event (Section~\ref{sec:bdi-agentspeak}).
\end{nbwhy}

\subsection{Terms, comparison, program database}
\label{sec:meta-programming}
\slidesrc{Prolog V: 1--24}

This deck covers the ``non-logical'' side of Prolog: predicates that inspect the
\emph{structure} of terms, compare them, test whether variables are bound, and read or modify
the \emph{program itself} while it runs.

\subsubsection{Handling terms in general: univ, functor, arg}
\slidesrc{Prolog V: 3--9}

Sometimes we want predicates that work on \emph{any} kind of term, without knowing its functor
in advance.

\begin{nbdef}{Univ \texttt{=..} (slide 3)}
\code{Term =.. List} (pronounced ``univ'') turns a term into a list: the first element is the
term's functor, the following elements are its arguments, in order. It also works in the
other direction.
\begin{minted}[linenos=false]{prolog}
?- p(1, q(3)) =.. X.       X = [p, 1, q(3)].
?- T =.. [p, 1, q(3)].     T = p(1, q(3)).
?- 1 =.. X.                X = [1].
?- [1,2,3] =.. X.          X = ['[|]', 1, [2,3]].      % slides: ['.', 1, [2,3]]
\end{minted}
\end{nbdef}

Since univ builds a list that is often not needed, it is usually better to use
\code{functor/3} and \code{arg/3}.

\begin{nbdef}{functor/3 (slides 4--5)}
\code{functor(Term, Name, Arity)} holds when the functor of \code{Term} is \code{Name} and its
arity is \code{Arity}. For atomic terms, the whole term is the name and the arity is 0. Either
\code{Term}, or both \code{Name} and \code{Arity}, must be bound, otherwise \code{functor}
throws an exception.
\begin{minted}[linenos=false]{prolog}
?- functor(p(1, q(3)), F, A).   F = p, A = 2.
?- functor(T, p, 2).            T = p(_, _).       % builds a "blank" term
?- functor(1, F, A).            F = 1, A = 0.
?- functor(X, 42, 0).           X = 42.
\end{minted}
\end{nbdef}

\begin{nbdef}{arg/3 (slides 6--7)}
\code{arg(N, Term, Arg)} holds if \code{Arg} is the \code{N}-th argument of \code{Term}
(counting from 1). \code{Term} must be bound to a compound term. In most systems \code{N} must
be bound; SWI-Prolog backtracks over all arguments. \code{arg/3} is deterministic when the
first two arguments are bound.
\begin{minted}[linenos=false]{prolog}
?- arg(2, foo(a,b), Arg).       Arg = b.
?- arg(N, foo(a,b), Arg).       N = 1, Arg = a ;  N = 2, Arg = b.
?- arg(N, 3, 3).                ERROR: Type error: `compound' expected, found `3'
\end{minted}
\end{nbdef}

\begin{nbexample}{collecting all functors of a ground term (slides 8--9)}
\begin{minted}{prolog}
all_functors(Term, List) :- all_functors(Term, List, []).

all_functors(Term, List, List0) :-
    functor(Term, Name, Arity),
    (   Arity > 0 ->
        List = [Name|List1],
        all_functors1(Arity, Term, List1, List0)
    ;   List = List0
    ).

all_functors1(N, Term, List, List0) :-
    (   N =:= 0 ->
        List = List0
    ;   arg(N, Term, Arg),
        all_functors(Arg, List1, List0),
        N1 is N - 1,
        all_functors1(N1, Term, List, List1)
    ).
\end{minted}
How to read it:
\begin{itemize}[nosep]
  \item \code{( Cond -> Then ; Else )} is Prolog's \textbf{if-then-else}: if \code{Cond}
        succeeds, run \code{Then}, otherwise \code{Else} (only the first solution of
        \code{Cond} is used);
  \item the pair \code{List}/\code{List0} is a \textbf{difference list}: ``the functors of
        this term, followed by \code{List0}''; it lets us concatenate without \code{append};
  \item \code{all\_functors1} visits the arguments from the last (\code{N} = arity) down to
        the first: the functors of argument \code{N} go in front of \code{List0}, and the
        functors of arguments $1 \ldots N-1$ go in front of those. So the final order is
        left to right.
\end{itemize}
\begin{minted}[linenos=false, fontsize=\footnotesize]{prolog}
?- all_functors(foo(bar([1,2],zip,baz(3)),bar(7)), L).
L = [foo, bar, '[|]', '[|]', baz, bar].         % slides: '.' instead of '[|]'
?- X = bar(42), all_functors(foo(X), L).
X = bar(42), L = [foo, bar].
?- all_functors(foo(X), L), X = bar(42).
ERROR: Arguments are not sufficiently instantiated
\end{minted}
In the first query: \code{foo}, then its first argument \code{bar(...)} gives \code{bar}, then
the list \code{[1,2]} = \code{'[|]'(1, '[|]'(2, []))} gives two list functors (numbers and
\code{[]} are atomic: none), \code{zip} is an atom (none), \code{baz(3)} gives \code{baz};
finally the second argument \code{bar(7)} gives \code{bar}. The last query fails because
\code{functor(X, Name, Arity)} is called with everything unbound: \textbf{the order of goals
matters} for these non-logical builtins.
\end{nbexample}

\subsubsection{Term comparison}
\slidesrc{Prolog V: 10}

\begin{nbfact}{Standard order of terms (slide 10)}
\code{@<}, \code{@=<}, \code{@>}, \code{@>=}, \code{==}, \code{\textbackslash==} work like
\code{<}, \code{=<}, \code{>}, \code{>=}, \code{=:=}, \code{=\textbackslash=} but on
\emph{any} terms (the sorting predicates are built on them):
\begin{itemize}[nosep]
  \item \textbf{Variables} \code{@<} \textbf{numbers} \code{@<} \textbf{atoms} \code{@<}
        \textbf{compound terms};
  \item atoms are compared lexicographically (alphabetically), numbers numerically;
  \item compound terms are compared first by \textbf{arity}, then by \textbf{functor name},
        then by the \textbf{arguments} in order;
  \item \code{==} and \code{\textbackslash==} check whether terms are \emph{identical},
        without binding any variables.
\end{itemize}
Beware of variables in terms being sorted!
\end{nbfact}

\begin{nbexample}{comparing terms}
\notinslides
\begin{center}
\small
\begin{tabularx}{\linewidth}{@{}l X@{}}
\toprule
\textbf{Query} & \textbf{Result and why} \\
\midrule
\code{X @< 1} & true (variables come first) \\
\code{42 @< abc} & true (numbers before atoms) \\
\code{apple @< banana} & true (alphabetical) \\
\code{z(1,1) @< a(1,1,1)} & true (arity 2 before arity 3, the name is not even looked at) \\
\code{f(1,b) @< f(1,c)} & true (same arity and name: compare arguments) \\
\code{X == Y} & false (two different unbound variables) \\
\code{X = Y, X == Y} & true (after unification they are the same variable) \\
\code{msort([b,2,f(x),a,1], L)} & \code{L = [1, 2, a, b, f(x)]} \\
\bottomrule
\end{tabularx}
\end{center}
\end{nbexample}

\subsubsection{Variable and type testing}
\slidesrc{Prolog V: 11--17}

\paragraph{The problem: \texttt{rev/2} in two directions.} The simple list reversal
\begin{minted}{prolog}
rev([], []).
rev([A|BC], R) :-
    rev(BC, CB),
    append(CB, [A], R).
\end{minted}
works ``forwards'' (\code{?- rev([1,2,3], R).} gives \code{R = [3,2,1]}), but used
``backwards'' (\code{?- rev(L, [3,2,1]).}) it goes into an \textbf{infinite loop after finding
the correct answer}: the recursive call \code{rev(BC, CB)}, with both arguments unknown,
produces longer and longer lists; once the right length has been found, no other length will
be correct, but Prolog keeps trying. Exchanging the calls:
\begin{minted}{prolog}
rev([], []).
rev([A|BC], R) :-
    append(CB, [A], R),
    rev(BC, CB).
\end{minted}
fixes the backwards mode (when \code{R} is bound, \code{append/3} has only one solution), but
now \code{rev/2} loops in the \emph{forwards} mode after finding the solution
(\code{append(CB, [A], R)} with \code{R} unbound generates longer and longer lists).

\begin{nbdef}{\texttt{var/1} and other non-logical builtins (slides 13--14)}
To implement a reversible version we must choose the goal order depending on \emph{which
arguments are bound}. \code{var(X)} succeeds if \code{X} is an unbound variable, and fails if
not: \code{?- var(\_).} is true, \code{?- var(a).} is false. Similar tests:
\begin{center}
\small
\begin{tabularx}{\linewidth}{@{}l X@{}}
\toprule
\code{nonvar(X)} & \code{X} is bound \\
\code{ground(X)} & \code{X} is bound, and every variable in its term is ground (no variables at all) \\
\code{atom(X)} & \code{X} is bound to an atom \\
\code{integer(X)}, \code{float(X)}, & \code{X} is bound to an integer / float / \\
\code{number(X)} & number \\
\code{atomic(X)} & \code{atom(X) ; number(X)} \\
\code{compound(X)} & \code{X} is bound to a compound term \\
\bottomrule
\end{tabularx}
\end{center}
All of them can make conjunction \textbf{not commutative}: \code{var(X), X = 1} succeeds,
\code{X = 1, var(X)} fails. Logically, ``A and B'' should not depend on the order.
\end{nbdef}

\begin{nbexample}{defining \texttt{length/2} in both directions (slides 15--17)}
When the list is supplied:
\begin{minted}{prolog}
len1([], 0).
len1([_|L], N) :-
    len1(L, N1),
    N is N1 + 1.
\end{minted}
When the length is supplied:
\begin{minted}{prolog}
len2([], 0).
len2([_|L], N) :-
    N > 0,
    N1 is N - 1,
    len2(L, N1).
\end{minted}
A version that works in both modes chooses with \code{var/1}:
\begin{minted}{prolog}
len(L, N) :-
    (   var(N) ->
        len1(L, N)
    ;   len2(L, N)
    ).
\end{minted}
NB: don't check whether \code{L} is a variable: even if it is not, its tail may be (e.g.\
\code{[a|T]}).
\begin{minted}[linenos=false]{prolog}
?- len(L, 4).          L = [_, _, _, _].
?- len([a,b,c], N).    N = 3.
?- len(L, N).          L = [], N = 0 ;  L = [_], N = 1 ;  L = [_, _], N = 2 ; ...
\end{minted}
Why the two helpers: \code{len1} cannot be used with \code{N} known and \code{L} unknown
(\code{N is N1 + 1} with \code{N1} unknown is an error, after the recursion produced
infinitely many candidates); \code{len2} needs \code{N} known (\code{N > 0} would raise an
error). The name \code{len} is used because \code{length/2} is a built-in that cannot be
redefined in SWI-Prolog.
\end{nbexample}

\subsubsection{Program access and manipulation}
\slidesrc{Prolog V: 18--24}

So far programs were assumed to be resident in memory, without asking how they are
represented or how they got there. Many applications need to \emph{access} the clauses of the
program, and to \emph{modify} it at run time (adding and deleting clauses).

\begin{nbdef}{clause/2, assert/1, asserta/1, retract/1 (slides 19--21)}
\begin{itemize}
  \item \code{clause(Head, Body)}: must be called with \code{Head} instantiated. The program
        is searched for the first clause whose head unifies with \code{Head}; its head and
        body are unified with \code{Head} and \code{Body}. On backtracking it succeeds once
        for each unifiable clause. Facts have the atom \code{true} as their body. Clauses
        cannot be accessed via their body.
  \item \code{assert(Clause)}: adds \code{Clause} as the \textbf{last} clause of the
        corresponding predicate (also \code{assertz}). \code{asserta(Clause)} adds it at the
        \textbf{beginning}. If \code{Clause} is uninstantiated, an error occurs.
  \item \code{retract(Clause)}: removes the first clause unifying with \code{Clause}. To
        retract \code{a :- b, c, d.} write \code{retract((a :- (b, c, d)))} or
        \code{retract((a :- Body))}.
\end{itemize}
\end{nbdef}

\begin{nbexample}{the program as data}
\notinslides
\begin{minted}[linenos=false, fontsize=\footnotesize]{prolog}
?- clause(grandparent(X, Y), Body).   Body = (parent(X, _A), parent(_A, Y)).
?- clause(parent(diana, C), Body).    C = william, Body = true ;  C = harry, Body = true.
?- assert(visited(rome)), assert(visited(paris)), asserta(visited(oslo)).
?- findall(C, visited(C), L).          L = [oslo, rome, paris].
?- retract(visited(rome)), findall(C, visited(C), L).   L = [oslo, paris].
\end{minted}
This is how a Prolog agent updates its beliefs when the world changes: \code{retract} the old
fact, \code{assert} the new one (AgentSpeak's \code{+b} and \code{-b},
Section~\ref{sec:bdi-agentspeak}).
\end{nbexample}

\begin{nbwarning}[Declare predicates you modify as dynamic]
In SWI-Prolog, predicates loaded from a file are \emph{static}: \code{assert}ing a clause of
\code{fib\_stupid/2} defined in the file raises ``No permission to modify static procedure''.
Add the directive \code{:- dynamic fib\_stupid/2.} at the top of the file (same for
\code{fib\_smart/2} and \code{visited/1}).
\end{nbwarning}

\paragraph{Assert in action: three Fibonacci programs (slides 22--24).} All three print a
message when the recursive clause is used, so we can count the calls. They have the
\code{N > 1} guard (Section~\ref{sec:arithmetic-nondeterminism}) after the \code{write}.
\begin{minted}{prolog}
fib_print(0, 1).
fib_print(1, 1).
fib_print(N, F) :-
    write('Call of fib_print on '), write(N), nl,
    N > 1,
    N1 is N - 1,
    N2 is N - 2,
    fib_print(N1, F1),
    fib_print(N2, F2),
    F is F1 + F2.
\end{minted}
\code{fib\_stupid} is identical but ends with \code{F is F1 + F2, assert(fib\_stupid(N, F))};
\code{fib\_smart} ends with \code{F is F1 + F2, asserta(fib\_smart(N, F))}.

\begin{nbexample}{what each version prints for \texttt{?- X(4, F).}}
\notinslides
\begin{center}
\small
\begin{tabularx}{\linewidth}{@{}l l X@{}}
\toprule
\textbf{Version} & \textbf{Printed calls} & \textbf{Why} \\
\midrule
\code{fib\_print} & 4, 3, 2, 2 & plain recursion: \code{fib(2)} is computed twice (once
  inside \code{fib(3)}, once as $N-2$ of 4). For larger $N$ the calls grow exponentially. \\
\code{fib\_stupid} & 4, 3, 2, 2 & \code{assert} adds the result \emph{after} the recursive
  clause: when \code{fib\_stupid(2, F)} is called again, the rule is tried \emph{before}
  the stored fact, so the memory is never used. A second query prints everything again
  (and adds duplicate facts). \\
\code{fib\_smart} & 4, 3, 2 & \code{asserta} puts \code{fib\_smart(2, 2)} \emph{first}:
  the second call \code{fib\_smart(2, F2)} matches the stored fact immediately. After the
  query, \code{fib\_smart(4, 5)}, \code{(3, 3)}, \code{(2, 2)} are stored: a second
  \code{?- fib\_smart(4, F).} prints nothing, and \code{fib\_smart(5, F)} prints only 5. \\
\bottomrule
\end{tabularx}
\end{center}
All three answer \code{F = 5}. \code{fib\_smart} is \textbf{memoisation} (dynamic
programming): each value is computed once, so the cost becomes linear.
\end{nbexample}

\subsection{Logical foundations}
\label{sec:logical-foundations}
\slidesrc{Prolog III: 1--9}

Why does Prolog compute correct answers? Because running a Prolog program is building a
\emph{proof} in first-order logic, with a single, very efficient inference rule:
\textbf{SLD resolution}. This deck was shown last; it ties Part~IV back to Part~II.

\subsubsection{Clauses and Horn clauses}
\slidesrc{Prolog III: 4--6}

\begin{nbdef}{Clause (slide 4)}
A \textbf{clause} is a disjunction of atoms or their negations (\textbf{literals}), where the
variables are implicitly \emph{universally quantified}:
\[
A_1 \lor A_2 \lor \cdots \lor A_n \lor \neg B_1 \lor \neg B_2 \lor \cdots \lor \neg B_m
\]
It is equivalent to
\[
(A_1 \lor \cdots \lor A_n) \leftarrow (B_1 \land \cdots \land B_m)
\]
usually written (conclusion $\leftarrow$ precondition) as
$A_1, \ldots, A_n \leftarrow B_1, \ldots, B_m$.
\end{nbdef}

Why are the two forms equivalent? With the rules of Section~\ref{sec:tautologies}:
\begin{align*}
&(A_1 \lor \cdots \lor A_n) \leftarrow (B_1 \land \cdots \land B_m) \\
&\quad\equiv \neg(B_1 \land \cdots \land B_m) \lor (A_1 \lor \cdots \lor A_n)
     && \reason{$q \leftarrow p \;\equiv\; \neg p \lor q$} \\
&\quad\equiv (\neg B_1 \lor \cdots \lor \neg B_m) \lor A_1 \lor \cdots \lor A_n
     && \reason{De Morgan} \\
&\quad\equiv A_1 \lor \cdots \lor A_n \lor \neg B_1 \lor \cdots \lor \neg B_m
     && \reason{commutativity}
\end{align*}
The positive literals end up on the left (conclusions), the negative ones on the right
(conditions).

\begin{nbdef}{Horn clause (slide 5)}
A \textbf{Horn clause} is a clause with \textbf{at most one positive literal} (a literal is a
predicate symbol applied to its arguments, or its negation).
\begin{center}
\small
\begin{tabular}{@{}l l l l@{}}
\toprule
\textbf{Kind} & \textbf{Logic} & \textbf{Prolog} & \textbf{Positive literals} \\
\midrule
definite clause: \textbf{rule} & $A \leftarrow B_1, \ldots, B_m$ & \code{a :- b1, ..., bm.} & 1 \\
definite clause: \textbf{fact} & $A \leftarrow$ & \code{a.} & 1 \\
negative clause: \textbf{goal} & $\leftarrow B_1, \ldots, B_m$ & \code{?- b1, ..., bm.} & 0 \\
\bottomrule
\end{tabular}
\end{center}
\end{nbdef}

\begin{nbfact}{Programs and questions (slide 6)}
\begin{itemize}[nosep]
  \item \textbf{Definite clauses represent knowledge}: a \textbf{logic program} is a set of
        definite clauses.
  \item \textbf{Negative clauses are questions}: the query \code{?- b1, ..., bm.} asks
        whether $\exists X_1 \ldots X_n.\,(B_1 \land \cdots \land B_m)$ holds (the
        $X_i$ are the variables of the query).
  \item To answer, Prolog uses \textbf{refutation}: it adds the \emph{negation} of the
        question and tries to derive a contradiction:
        \begin{align*}
        &\neg\,\exists X_1 \ldots X_n.\,(B_1 \land \cdots \land B_m) \\
        &\quad\equiv \forall X_1 \ldots X_n.\,(\neg B_1 \lor \cdots \lor \neg B_m)
          && \reason{duality, De Morgan} \\
        &\quad\equiv\; \leftarrow B_1, \ldots, B_m
          && \reason{a negative clause}
        \end{align*}
\end{itemize}
\end{nbfact}

\begin{nbintuition}[Refutation = proof by contradiction]
It is the $\neg$\,Intro idea of natural deduction (Section~\ref{sec:natural-deduction-prop}):
to show that ``someone is mortal'', assume ``nobody is mortal'' and derive a contradiction
from it and the program. The contradiction is the \textbf{empty clause} $\Box$ (a disjunction
with no literals, which can never be true). The bindings used along the way tell \emph{who}
is mortal: they are the answer.
\end{nbintuition}

\subsubsection{SLD resolution}
\slidesrc{Prolog III: 7}

\begin{nbdef}{SLD resolution (slide 7)}
\textbf{S}election-function, \textbf{L}inear resolution for \textbf{D}efinite clauses.
Start from a definite program $P$ and a goal
\[ G_0 = \;\leftarrow A_1, \ldots, A_m \qquad (m \geq 0). \]
\begin{enumerate}[nosep]
  \item The \textbf{selection rule} selects a sub-goal $A_i$.
  \item Choose a program clause, \textbf{renamed} with fresh variables,
        $B_0 \leftarrow B_1, \ldots, B_n$ ($n \geq 0$), whose head $B_0$ unifies with $A_i$,
        with most general unifier $\theta_1$.
  \item The new goal is
        \[ G_1 = \;\leftarrow (A_1, \ldots, A_{i-1}, B_1, \ldots, B_n, A_{i+1}, \ldots, A_m)\,\theta_1 \]
        i.e.\ $A_i$ is \emph{replaced by the body} of the clause, and the unifier is applied
        to the whole goal.
\end{enumerate}
Repeat until the goal is empty ($\Box$: success) or no clause unifies with the selected atom
(failure of that branch). The answer is the composition of all the unifiers, restricted to
the variables of $G_0$.
\end{nbdef}

\begin{nbexample}{SLD derivation of \texttt{?- mortal(X).}}
Program: \code{mortal(X) :- human(X).} \ \code{human(socrates).}
\begin{align*}
G_0 &= \;\leftarrow \mathit{mortal}(X) \\
  &\quad \text{clause renamed } \mathit{mortal}(X_1) \leftarrow \mathit{human}(X_1),\quad
     \theta_1 = \{X_1 / X\} \\
G_1 &= \;\leftarrow \mathit{human}(X) \\
  &\quad \text{clause } \mathit{human}(\mathit{socrates}) \leftarrow,\quad
     \theta_2 = \{X / \mathit{socrates}\} \\
G_2 &= \;\Box \qquad \text{(empty goal: success)}
\end{align*}
Answer: $X = \mathit{socrates}$. It is Whiteboard~\ref{wb:modus-ponens} again: unification
does the instantiation, resolution does the Modus Ponens.
\end{nbexample}

\subsubsection{SLD trees}
\slidesrc{Prolog III: 8--9}

\begin{nbdef}{SLD tree (slide 8)}
Let $P$ be a definite program, $G_0$ a goal and $\mathcal{R}$ a selection rule. The SLD tree
of $G_0$ (using $P$ and $\mathcal{R}$) is a (possibly infinite) labelled tree such that:
\begin{itemize}[nosep]
  \item the root is labelled by $G_0$;
  \item if $\leftarrow A_1, \ldots, A_n$ is a node and $A_m$ is the atom selected by
        $\mathcal{R}$, the node has one child for each clause $A \leftarrow B_1, \ldots, B_q$
        of $P$ such that $A$ and $A_m$ are unifiable; the child is
        \[ \leftarrow (A_1, \ldots, A_{m-1}, B_1, \ldots, B_q, A_{m+1}, \ldots, A_n)\,\theta \]
        where $\theta$ is the mgu of $A_m$ and $A$.
\end{itemize}
Leaves labelled $\Box$ are \textbf{success} leaves (each gives an answer); leaves where no
clause applies are \textbf{failure} leaves.
\end{nbdef}

(The slide writes the child as $(A_1, \ldots, A_{m-1}, B_1, \ldots, A_{m+1}, \ldots, A_n)\theta$,
dropping $B_q$: the whole body $B_1, \ldots, B_q$ replaces $A_m$.)

\begin{nbexample}{SLD tree with the \emph{rightmost} selection rule (slide 9)}
Program (colours as on the slide):
\begin{center}
\textcolor{nbred}{(1)\ $p(X, Z) \leftarrow p(Y, Z), q(X, Y)$} \qquad
\textcolor{nbblue}{(2)\ $p(X, X)$} \qquad
\textcolor{nbgreen!60!black}{(3)\ $q(a, b)$}
\end{center}
Goal $\leftarrow p(X, b)$; the selection function picks the \textbf{rightmost} atom.
\begin{center}
\begin{forest}
for tree={draw, rounded corners, font=\small, l sep=8mm, s sep=10mm, edge={thick}, align=center}
[{$\leftarrow p(X, b)$}, fill=nbblue!8
  [{$\leftarrow p(Y, b),\ q(X, Y)$}, fill=nbblue!8, edge={nbred}, edge label={node[midway, left, font=\scriptsize]{(1)}}
    [{$\leftarrow p(b, b)$}, fill=nbblue!8, edge={nbgreen!60!black}, edge label={node[midway, right, font=\scriptsize]{(3): $X = a,\ Y = b$}}
      [{$\leftarrow p(Y', b),\ q(b, Y')$\\failure: no $q(b, \_)$}, fill=nbred!10, edge={nbred}, edge label={node[midway, left, font=\scriptsize]{(1)}}]
      [{$\Box$\\answer $X = a$}, fill=nbgreen!15, edge={nbblue}, edge label={node[midway, right, font=\scriptsize]{(2)}}]]]
  [{$\Box$\\answer $X = b$}, fill=nbgreen!15, edge={nbblue}, edge label={node[midway, right, font=\scriptsize]{(2): $X = b$}}]]
\end{forest}
\end{center}
Step by step:
\begin{enumerate}[nosep]
  \item Root $\leftarrow p(X, b)$: only one atom. Clause (1), renamed
        $p(X_1, Z_1) \leftarrow p(Y_1, Z_1), q(X_1, Y_1)$, unifies with $X_1 = X$, $Z_1 = b$:
        child $\leftarrow p(Y, b), q(X, Y)$ (writing $Y$ for $Y_1$). Clause (2), $p(X_2,
        X_2)$, unifies with $X_2 = X = b$: child $\Box$, answer $X = b$.
  \item Node $\leftarrow p(Y, b), q(X, Y)$: rightmost atom $q(X, Y)$; only clause (3)
        unifies: $X = a$, $Y = b$; child $\leftarrow p(b, b)$.
  \item Node $\leftarrow p(b, b)$: clause (1) gives $\leftarrow p(Y', b), q(b, Y')$, whose
        rightmost atom $q(b, Y')$ unifies with no clause: failure. Clause (2) unifies
        $p(b, b)$ with $p(X_2, X_2)$: $\Box$, answer $X = a$.
\end{enumerate}
The tree is \emph{finite}: two answers, $X = b$ and $X = a$.
\end{nbexample}

\begin{nbwarning}[What real Prolog does with the same program]
\notinslides
Prolog's selection rule is \textbf{leftmost}, its search is \textbf{depth-first} with clauses
tried \textbf{top to bottom}. With the leftmost rule the first child of
$\leftarrow p(X, b)$ is $\leftarrow p(Y, b), q(X, Y)$, whose leftmost atom $p(Y, b)$ is again
resolved first with clause (1): $\leftarrow p(Y', b), q(Y, Y'), q(X, Y)$, and so on. The
leftmost branch of the SLD tree is \emph{infinite}, and depth-first search explores it first:
\code{?- p(X, b).} never returns an answer (stack overflow), although $X = b$ and $X = a$ are
logical consequences of the program.
\begin{itemize}[nosep]
  \item The \emph{set of answers} does not depend on the selection rule (every SLD tree
        contains the same success leaves), but the \emph{shape} of the tree, and hence
        termination, does.
  \item Depth-first search is fast and memory-cheap, but \emph{incomplete}: it can get lost in
        an infinite branch before reaching a success leaf. This is the ``left recursion''
        problem of Section~\ref{sec:prolog-basics}. Writing clause (2) first, or the body as
        \code{q(X, Y), p(Y, Z)}, would make Prolog find both answers.
\end{itemize}
\end{nbwarning}

\begin{center}
\begin{tikzpicture}[font=\small, >=Stealth,
  b/.style={draw, thick, rounded corners, text width=3.6cm, align=center, minimum height=1.3cm, fill=#1}]
  \node[b=nbblue!8] (fol) at (0,0) {\textbf{FOL}\\full first-order logic, undecidable};
  \node[b=nbpurple!8] (horn) at (4.5,0) {\textbf{Horn clauses}\\facts, rules, goals};
  \node[b=nbgreen!10] (sld) at (9,0) {\textbf{SLD resolution}\\+ leftmost, depth-first\\= \textbf{Prolog}};
  \draw[->, thick] (fol) -- node[above, font=\scriptsize]{restrict} (horn);
  \draw[->, thick] (horn) -- node[above, font=\scriptsize]{one rule} (sld);
\end{tikzpicture}
\captionof{figure}{From logic to Prolog: restrict FOL to Horn clauses, use one inference rule
(SLD resolution), and fix a search strategy. The price of efficiency is that some logically
valid answers may not be found (non-termination).}
\end{center}


\begin{nbsummary}[Prolog]
\small
\begin{itemize}[nosep]
  \item Program = \textbf{facts} (\code{parent(diana, harry).}) + \textbf{rules}
        (\code{H :- B1, B2.}, ``H if B1 and B2''). Query \code{?- G.}: Prolog tries to prove
        \code{G}, returning variable bindings; \code{;} asks for more answers.
  \item \textbf{Variables} start with upper case or \code{\_}; they are unknowns,
        \emph{bound once} (no assignment). \code{\_} is anonymous. Arguments can be input or
        output (``multiple directions'').
  \item \textbf{Terms}: atoms, numbers, variables, compound terms \code{f(t1,...,tn)}
        (functor/arity, written \code{name/arity}). Lists: \code{[]}, \code{[H|T]}.
  \item \textbf{Unification} \code{=}: makes two terms identical by binding variables;
        used for argument passing. \code{==} tests identity without binding.
        \code{is} evaluates arithmetic; \code{=:=}, \code{<}, \code{>=}\ldots{} compare numbers.
  \item Execution: goals left to right, clauses top to bottom, depth-first search with
        \textbf{backtracking}. Order of goals/clauses affects efficiency and
        \emph{termination}, not the logical meaning.
  \item Several clauses are \textbf{not} an if-then-else: every applicable clause can give
        an answer (absolute value bug); add guards (\code{X < 0}).
  \item \code{member/2}, \code{append/3}, \code{length/2}; accumulators for efficiency;
        \code{findall/3}, \code{bagof/3}, \code{setof/3} collect all solutions.
  \item Meta: \code{=..} (univ), \code{functor/3}, \code{arg/3}; standard order
        \code{@<}; type tests \code{var}, \code{nonvar}, \code{atom}, \ldots;
        \code{clause/2}, \code{assert/asserta}, \code{retract} (declare \code{:- dynamic}).
  \item Foundations: a clause is a disjunction of literals; \textbf{Horn clause}: at most
        one positive literal (facts, rules, goals). \textbf{SLD resolution}: select an atom
        of the goal, resolve it with a renamed program clause using the mgu; an
        \textbf{SLD tree} shows all choices; success leaves $\Box$ give answers. Prolog =
        leftmost selection + depth-first, top-to-bottom search.
\end{itemize}
\end{nbsummary}
